\section{Orientations}

\paragraph{Point 1}
\begin{quotex}
Tonight, we are posting the first point of \textbf{Julius Evola}'s Orientations and will continue for the next ten Sundays. 

\end{quotex}
It is useless to create illusions with the pipe dream of any optimism whatsoever: we find ourselves today at the end of a cycle. Already for centuries, at first unfelt, then with the surge of a mass that caves in, numerous processes in the West have destroyed every normal and legitimate ordering of men, and distorted every higher conception of living, acting, knowing, and struggling. And the impulse of this downfall, its speed and its dizziness, is called ``progress." And hymns were sung to ``progress" and we were deluded that this civilization is — a civilization of matter and machines — civilization par excellence. The entire history of the world was preordained to it: until the ultimate consequences of this whole process was such as to bring a few to an awakening.

Where, and under what symbols, of those who sought to organize the forces for a possible resistance, is well-known. On the one hand, there is a nation [Italy] which, from when it became unified, had known only the mediocre climate of liberalism, democracy, and constitutional monarchy — dared take up again the symbol of Rome [the eagle] as the basis for a new political conception and for a new ideal of virility and dignity. Analogous forces were awakened in the nation [Germany] that had made in the Middle Ages the Roman symbol of the Imperium as its own, to reassert the principle of authority and the supremacy of those values that have their root in blood, in race, and in the deepest strength of a people. And while in other European nations some groups were already oriented in the same direction, a third force in the Asian continent [Japan] added itself to the alliance, the nation of samurai, in which the adoption of exterior forms of modern civilization did not compromise the fidelity to a warrior tradition centered on the symbol of the solar Empire of divine right.

We do not pretend that in these currents the distinction between the essential and the incidental is clear cut, that in them it made to the ideas of the opposing party an adequate persuasion and qualification of persons, that there were various outdated influences affected by the same forces that had to be combated. The process of ideological purification would have been able to take place a second time, so that some immediate and non-deferrable political problems were solved. But it was also clear that an alliance of forces was taking form, representing an open challenge to ``modern" civilization: both to that of the democracies inherited from the French Revolution and to the other, representing the extreme limit of the degradation of Western man: the collectivistic civilization of the Fourth Estate, the communist civilization of the mass-man without a face.

The rhythms accelerated and the tensions increased ending up in an armed conflict of these forces. What prevailed was the massive power of a coalition that did not back down in the face of the most hybrid of their agreements, and of the most hypocritical of this ideological mobilization in order to crush the world that was rising up and that intended to assert its right. Whether or not our men were equal to the task, if errors were committed in regard to timeliness, complete preparation, or measure of risk, that is beside the point, that is not something that compromised the inner meaning of the battle that was fought. Likewise, it does not interest us that today history is avenged on the victors, that the democratic powers which, having formed a coalition with the forces of red subversion just to conduct the war up to the senseless extremism of an unconditional surrender and total destruction, today sees come back to roost against the allies of yesterday, a danger more terrible than the one they wished to avoid.

The only thing that counts is this: we find ourselves today in the middle of a world of ruins. And the problem to be posed is: do men still exist on their feet within the middle of these ruins? And what should they do, what else can they still do?

\paragraph{Point 2}

\begin{quotex}
The way out of the current situation, according to Evola, will not come from a new political party or platform. Rather, it will require men of a certain character, vision of life, and principles. These are not yet made. Probably, they are still not clear, because after 60 years, there is still no single vision of life among those who oppose the modern world.

Political systems alone are not determinate: the best systems will fail if the population cannot maintain it, while a superior people will succeed despite the system. After 60 years, there should be evidence for this claim. In making this moral argument, Evola goes against the entire modern world which sees the Axis powers he supported as the very embodiment of evil. 

\end{quotex}
Such a problem indeed goes beyond yesterday's alliances, since it is clear that the winners and losers [\emph{vincitori e vinti}] now find each other on the same level and that the only result of the second world war was the reduction of Europe to the object of extra-European powers and interests. It must be recognized then that the devastation that we have around us is of an especially moral character. We are in a climate of general moral anesthesia, of profound disorientation, in spite of all the words of order in use in a society of consumption and democracy: the breakdown of character and every true dignity, the ideological confusion, the prevalence of the lowest interests, living hand to mouth, are what characterizes, in general, the post-war man. To recognize this, also means to recognize that the first problem, the base of every other, is of an inner nature: To pick oneself up, to rise up interiorly, to give oneself a form, to create in oneself order and rectitude.

No one has learned from the lessons of the recent past if he still fools himself about the possibilities of a purely political battle and about the power of one or the other formula or system, whose new human quality is not made by a specific opposing party. Here is a principle that today more than ever should have absolute evidence: if a State possesses a political or social system that, in theory, is of value as the most perfect, but the human element was deficient, that state would descend sooner or later to the level of the basest societies. However, a people, a race capable of producing true men, men of right feeling and reliable instinct, would reach a high level of civilization and would stand on their feet in the face of the most calamitous tests even if its political system was defective and imperfect. One takes therefore a precise position against that false ``political realism" that thinks only in terms of programs, party organizational problems, or social and economic prescriptions. All this belongs to the contingent, not to the essential. The measure of what can still be saved depends upon the existence — or not — of men who are in front not to preach formulas, but to be examples, not going towards demagogy and the materialism of the masses, but to awaken different forms of sensibility and interests. Starting from what can still remain among the ruins, to reconstruct slowly a new man to animate through a determinate spirit and a suitable vision of life, to fortify through the tenacious adherence to given principles — this is the true problem.

\paragraph{Point 3}
\begin{quotex}
First of all, what is that ``legionary spirit" and where is the evidence for it today? That spirit has nothing to do with hot-headedness, emotional outbursts, or partisanship. Rather, it is characterized by an iron will, a sense of order, rationality, mental sobriety, the ability to plan or strategize, loyalty, the adherence to principles.

How does that fit in today, among those who can see in such a spirit only a disease, a mistake, the ``authoritarian personality". Today hedonism is preferred to sobriety, living for the moment instead of strategizing, the lack of discrimination instead of the quest for order, the irrational, the unusual, the absurd instead of rationality. But what idea or principle is there to be loyal to? Even among certain groups that claim to be ``radical traditionalists" there is no underlying unity. Rather, they are motivated by the party spirit, i.e., alliances are made based on similarity of immediate goals, but not on the possession of a common mind.

Evola's vision is far from being realized, even further than when Orientations was published. Where do we see the \emph{legionary spirit} ``in the factories, in the laboratories, in the universities, in the street"? The ``agitators" today wonder why no one will follow them. Isn't it time to look elsewhere? 

\end{quotex}
Something already exists as spirit that can serve as evidence of the forces of resistance and resurgence: it is the \emph{legionary spirit}. It is the attitude of those who were able to choose the hardest way, who were able to fight even knowing that the battle was materially lost, who were able to confirm the words from the ancient saga: ``Loyalty is stronger than fire" and through which the traditional idea was affirmed; that is the meaning of honor and shame, — not small measures drawn from petty morality — that creates a substantial, existential difference between beings, almost like that between one race and another.

On the other hand, there is the realization typical of those in whom what was the end now appeared as the means, in them the recognition of the illusory character of multiple myths while leaving untouched those who were able to achieve \emph{for themselves}, on the borderline between life and death, beyond the world of contingency.

These forms of the spirit can be the base of a new unity. The essential thing is to adopt them, to apply them, and to extend them from the time of war to the time of peace, of this peace especially, that is only a setback and a poorly kept down disorder—to those who bring about distinctions and a new alignment. That must happen in rather more essential terms than what is a ``party", which can only be a contingent instrument in sight of given political battles; in terms more essential than even as a simple ``movement", if by ``movement" is meant only a phenomenon of the masses and aggregation, a quantitative phenomenon more than qualitative, based more on emotive factors than on strict, clear adherence to an idea.

It is instead a silent revolution proceeding in the depths, that must be propitiated, to those who are, it first creates on the inside and in the individual the premises of that order that then will have to also be asserted on the outside, supplanting as quick as lightning at the right moment the forms and forces of a world of subversion. The ``style" that must gain emphasis is that of whoever holds onto positions in faithfulness to himself and to an idea, in a collected intensity, in a repulsion for every compromise, in a total engagement that must be manifested not only in the political battle, but also in every expression of existence: in the factories, in the laboratories, in the universities, in the street, in the very personal life of the affected. One must reach the point that the type of which we speak, and who must be the cellular substance of our alliance, is well recognizable, unmistakable, differentiated, and one can say: ``He is one who acts like a man of the movement."

This must today be recovered, which was already the guarantee of the forces that yearned for a new order for Europe but which in its realization often was impeded and diverted by multiple factors. And today, fundamentally, the conditions are better, because misunderstandings do not exist and it suffices to look around, from the streets to the Parliament, because the callings are put to the test and one has, clearly, the measure of what we must \emph{not} be. Facing a world of mush whose principle is: ``Who makes you do it", or: ``First comes the stomach, the skin (the \emph{Malapartian Skin}!\footnote{\url{http://www.britannica.com/EBchecked/topic/359524/Curzio-Malaparte}}), and then morality", or again: ``These are not times in which we can be permitted the luxury of having a character", or finally: ``I have a family". One is able to oppose clarity and service to them: ``We cannot do otherwise, this is our life, this is our being". Whatever is positive will be able to be reached today or tomorrow. It will not be through the abilities of agitators and politicians, but rather through the natural prestige and the recognition of men whether of yesterday, or, and even more, of the new generation, that so many are capable and in that give the guarantee for their idea.

\paragraph{Point 4}
\begin{quotex}
Here a new vision is put before our eyes. The regression of castes must be walked back up, but in today's undifferentiated world, one's caste is unclear. New spirits will awaken, but not from the expected places. This has nothing to do with economic classes. It is a revolt against the conformity of the bourgeois world, but from superior values, not the lower values of the mass man.

This awakening spirit comes from the inside and cannot be identified with an external movements. Thus all attempts to identify such spirits based on adherence to some political ideology or religious form will be far off the mark. Yet, these spirits will recognize each other and form natural hierarchic structures. It is not a ``big tent" or equals, or debaters in endless conversation. Hence, counter-traditional movements that try to include anarchists or the abnormal cannot be part of this.

Sentimentality has no part of this, despite those who tell me I should ``like" this person or web site, etc. But ``liking" is not the issue. This sentimentality is opposed to ``impersonality", the idea that it is ``strictly business, nothing personal"; this attitude comes apparently from a deep part of Evola and its source is not the Nordic soul.

The text follows. 

\end{quotex}
Therefore, a new substance must clear the way in slow advance beyond the squares, the classes, and social positions of the past. It is necessary to keep a new pattern right in front of your eyes, in order to measure your own strength and calling. It is important and fundamental to recognize exactly that this pattern has nothing to do with classes as economic categories, and the antagonisms related to them. It could manifest itself in the clothing of the rich as well as the poor, of the worker as well as the aristocrat, of the businessman as well as the explorer, of the technician, theologian, farmer, or the politician in the strict sense. But this new substance will know an internal differentiation, which will be perfect when, again, there will be no doubt about the vocations and functions of followers and commanders, when a restored symbol of unshaken authority will tower over the center of new hierarchic structures.

That defines a direction that is as anti-bourgeois as anti-proletarian, a direction totally devoid of democratic contaminations and ``social" foibles, because leading toward a clear, virile, articulated world, made by men and leaders of men. Disdain for the bourgeois myth of ``security", for the petty standardized, conformist, domesticated, and ``moralized" life. Disdain for the anodyne bond characteristic of every collectivistic and mechanistic system and of all ideologies that grant to confused ``social" values primacy over those heroic and spiritual values by which we have to define, in every domain, the type of true man, of the absolute person.

And something essential will be attained when the love for a style of active impersonality is revived, for which that which counts is the work and not the individual, for which a man is not capable of considering himself as something of importance, importance being instead his function, his responsibility, the task he adopts, the goal he pursues. Where this spirit asserts itself, many problems of the economic and social order will be simplified, which would otherwise remain insoluble if confronted from the outside, without the counter party of a change of spiritual factors and without the elimination of infectious ideologies that already in part prejudice every return to normality, as well as the perception itself of what normality means.

\paragraph{Point 5}
\begin{quotex}
Although the Cold War between the USSR and USA is over, this makes Evola's insights in this point all the more prescient. As long as communism and the Soviet Union were perceived as external enemies, the USA could sustain internally a certain amount of reaction and conservatism to define itself as the opposite of the USSR. However, once the wall came down, the right in the USA became irrelevant, so that progressive forces could continue unchecked, no longer subject to the charge of being pro-marxist and treasonous.

Evola was not deceived and recognized that the various movements, ostensibly in opposition, were in fact different degrees of the same frame of mind. The imagery he uses of the dusk and poison are beautifully written. Little else has changed in 60 years, since the trends he abhors are still considered to be ``progress", and the West is in a constant state of ``waiting" … waiting for the culmination of this ``progress" in the indefinite future in who knows what form. Since the ``real" world only exists in the future, the past, and Tradition, are thereby disvalued. Evola may be placing too much blame on the USA, for Europe herself has been equally complicit, if not even leading the way. He is correct, though, in a significant point: the battle has been lost without even a shot being fired. 

\end{quotex}
Not only as a doctrinal orientation, but also in regard to the world of action, it is then important for men of the new arrangement to recognize correctly the concatenation of causes and effects and the essential continuity of the trend that has given life to the various political forms in play today in the chaos of parties. Liberalism, then democracy, then socialism, than radicalism, finally communism and Bolshevism have historically appeared only as degrees of the same evil, as stages that prepare each one that follows in the complex of a process of decline. And the beginning of this process stands at the point at which Western man broke his ties to tradition, denying every higher symbol of authority and sovereignty; he claimed for himself as an individual a vain and illusory freedom; he becomes an atom rather than a conscious part of the organic unity and hierarchy of a whole. And the atom, in the end, had to find, opposed to himself, the mass of other atoms, other individuals, and to be involved in the emergence of the reign of quantity, of the pure number, of the materialized masses and not having another God outside of the sovereign economy.

This process does not stop here halfway. Without the French Revolution and liberalism, there would not have been constitutionalism and democracy, without democracy, there would not have been socialism and demagogic nationalism, without the preparation of socialism there would not have been radicalism and finally communism. The fact that these various forms today often represent, one alongside the other or in opposition, must not prevent the recognition, to an eye that truly sees, that they hold themselves together, they link to each other, mutually condition each other and express only the different degrees of the same trend, the same subversion of every normal and legitimate social order.

So the great illusion of our time is that democracy and liberalism are the antitheses of communism and have the power to stem the tide of the forces from below, of what in the jargon of the syndicates are called the ``progressive" movement. Illusions: like those who say that the dusk is the opposite of the night, that the incipient degree of an evil is the antithesis of the acute and endemic form of it, that a diluted poison is the antidote of the same poison in its pure and concentrated state. The men in the government of this ``liberated" Italy have learned nothing from recent history whose lessons were repeated everywhere to the point of monotony, and continue their pitiful game with out-of-date political conceptions and hands in the parliamentary carnival, almost a dance macabre on a simmering volcano.

But for \emph{us} there must instead be the courage of radicalism, the ``no" said to political decadence in all its forms, both of the left and of a presumed right. And, above all, we must be aware of that which he does not share with subversion, and that making concessions today means condemning oneself to be totally overwhelmed tomorrow. The rigour of the idea, therefore, and readiness to push ahead by pure force when the right moment is reached.

That naturally also implies the disengagement from ideological distortions, unfortunately popular even in a part of the youth, because of which they indulge in excuses for the destruction that has already happened, deluding themselves by thinking that they, after all, were necessary and served ``progress"; they think that it must be combated for anything ``new", put off into a determined future, rather than for the truth that we already possess because they, although in various forms of application, always and everywhere made from the base of every right type of social and political organization.

These foibles must be rejected. And we come back to those who are accused of being ``anti-historical" and ``reactionary". History does not exist, a mysterious entity spelled with a capital H. They are the men, if even they are \emph{really} men, who make and unmake history; so called ``historicism" is more or less the same thing as what is called ``progressivism" in leftist circles and it wants a single thing today: to foment passivity in respect to the trend that swells and always leads lower. And, as to ``reactionism", you ask: You would therefore prefer that while you act, destroying and profaning, we do not ``react", but stay to look so that they say to you: bravo!, continue? We are not ``reactionary" only because the word is not strong enough especially because we start from the positive, we represent the positive, real and original values, the light of some ``sun of the future" is not necessary.

In the face of our radicalism, in particular, the antithesis between the red ``East" and the democratic ``West" appears irrelevant, yet a possible armed conflict between these two blocks appears to us also tragically irrelevant it. To look only at the immediate, the choice of the lesser evil certainly remains because the military victory of the ``East" would imply the immediate physical destruction of the last exponents of resistance. But at the base of the idea, Russia and North America are to be considered as two prongs of the same pincer because they squeeze themselves permanently around Europe. In two different but converging forms, the same foreign and hostile force acts in them. The forms of standardization, conformism, democratic leveling, frenetic productivity, a more or less domineering and explicit ``brain trust", and piecemeal materialism in Americanism can only serve to pave the road for the last phase that is represented, along the same direction, by the communist ideal of the mass-man.

The distinctive character of Americanism is that the attack against quality and personality is not put into effect by the brute force of a Marxist dictator and the mind of the State, but it happens almost spontaneously, by the ways of a civilization ignorant of ideals higher than wealth, consumption, productivity, production without brakes. Therefore, through a worsening and a reduction to the absurd of what Europe herself chose, the same patterns have taken, or are taking, form there. But primitivism, mechanism, and brutality remain so much more on one side than the other. In a certain sense, Americanism is more dangerous for us than communism, because it is a type of Trojan horse. When the attack against the residual values of European tradition is effectuated in the direct and naked forms characteristic of Bolshevik ideology or Stalinism, some reactions still reawaken, certain lines of resistance, even if ephemeral, can be maintained. Things are different when the same evil acts in a more subtle way and the transformations take place on the plane of the general customs and visions of life, as is the case for Americanism. Enduring this influence with a light heart in the name of democracy, Europe is already predisposed for the final abdication, so that it could even happen without even the need for a military disaster, but it will arrive, through ``progressive" path, after a final social crisis, at more or less at the same point. Again, it doesn't stop here halfway. Americanism, willing it or not, works for its apparent enemy, for collectivism.

\paragraph{Point 6}
\begin{quotex}
There are several important issues brought in Point 6. That these are in dispute today is incredible, if, in fact, they are even noticed. Perhaps the professional politicians do know it, but are coy enough not to mention them publicly.

\begin{itemize}
\item
The first is the reduction of the political to the economical. This is even more pronounced in the USA elections, where the economy is supposed to outweigh every other issue. The net effect of this is that, because they are left out of the political discussion, lower political and social forces can continue unchecked. 

\item
The second is that there is no such thing as a ``spontaneous" grassroots movement from below. We are led to believe that there are certain important issues and movements then arise to deal with them. The opposite is actually the case; movements fabricate issues to gain power and achieve their own ends. 

\item
The false ideal of ``social justice" is still used to this day for political gain. It is quite effective, except for those who expect ``justice". 

\item
On a positive note, Evola proposes, as always, a hierarchical arrangement based on high and heroic ideals. To replace the parliamentary system of professional politicians, he proposes a form of corporativism. This is not, as often believe, the rule by corporations as understood in the West, particularly the USA. Rather, it refers to the representation of the various business sectors, e.g., health, energy, etc., by workers as well as owners. They would bring expertise in their fields to the political and spiritual hierarchy, unlike democratic politicians who typically have no expertise. Contra this, one could complain that these are ``special interests" who already have too much influence. However, now it is done in secret, behind closed doors. That would not be the case in a corporative system which actually has a more balanced and diverse representation than current Western deliberative bodies. Furthermore, in such a system, economic factors alone would not have dominance. 
\end{itemize}
\end{quotex}
Not without relation to that, our radicalism of the reconstruction requires that we do not compromise not only with every variety of Marxist or socialist ideology, but also with what in general can be called the \emph{hallucination} or the \emph{demon of the economy}. Here it is a question of the idea that in both individual and collective life, the economic factor is the important, real, and decisive one; that the concentration of every value and interest on the economic and productive plane is not the unprecedented aberration of modern Western man, but rather something normal, not a possible brute necessity, but something that must be desired and exalted. Both capitalism and marxism remain enclosed in this dark circle. We must break through this circle. As long as we know how to speak only of economic classes, work, salary, production, as long as we delude ourselves that true human progress and the true elevation of the individual are conditioned by a particular system of distribution of wealth and goods and have therefore something to do with indigence or affluence, with the state of USA \emph{prosperity,} or with that of utopian socialism, we always remain on the same level of what must be combated. This we must affirm: everything that is economy and economic interest as mere satisfaction of physical needs had, has, and always will have a subordinated function in normal humanity; that beyond this sphere an order of higher, political, spiritual and heroic values must be differentiated, an order that–as we already said—does not know and not even admits, proletarian or capitalists, and only depends on what things must be defined as worth living and dying for. A true hierarchy must be established, new dignities must be differentiated and, at the top, a higher function of command, of the \emph{imperium} must dominate.

Thus, in such regards, many evil herbs that have taken root here and there are eradicated, sometimes even in our camp. What is, in fact, this talk of ``State of work", of ``national socialism", of ``humanism of work" and similar ideas? What are these more or less stated requests for a regression of politics into the economy, almost as a resumption of those problematic tendencies toward an integral and, fundamentally, headless, corporatism that fortunately already found the way blocked in fascism? What is this consideration of the formula of ``socialization" as a type of universal drug and this elevation of the ``social idea" to the symbol of a new civilization that, who knows how, should be beyond both ``East" and ``West"?

These—it is necessary to recognize this—are the sides of a shadow present in not a few spirits, who only, for other reasons, find themselves on the same side as us. With that they think they are faithful to a revolutionary order, while they obey only stronger suggestions of them with which a degraded political environment is saturated. And among such suggestions the same ``social question" returns. When will they finally take into account the truth, i.e., that Marxism did not arise because a real social question existed, but the social question rises—in endless cases—only because Marxism exists, that is to say, artificially, and yet in almost always insoluble terms, to the work of agitators, to the famous ``awakeners of class consciousness", on which Lenin expressed himself very clearly, when he refuted the spontaneous character of proletarian revolutionary movements?

It is starting from this premise that it would be necessary to act, in the first and foremost meaning of ideological anti-proletarianization, as the disinfection of the still healthy parts of the people from the socialist virus. Only then will the one or the other reform be able to be studied and actuated without danger, in compliance with true justice.

So, as the particular case, we will see, according to that spirit, the corporative idea can be again one of the bases of the reconstruction: corporativism not so much as a general system of state and almost bureaucratic composition that maintain the deleterious idea of opposed political arrangements of class systems, but rather as the need that in the very interior of the business that unity, that solidarity of differentiated forces be reconstructed, that the capitalist lie (with the subversive parasitic type of the speculator and the finance-capitalist) on one side, the Marxist agitation on the other, have jeopardized and shattered. It is necessary to bring the business to the form of an almost military unity, in which they compare the solidarity and the fidelity of associated working forces around it in the common enterprise to the spirit of responsibility, to the energy and the competence of the directors. The only true task is, however, the \emph{organic reconstruction of the business}, and to realize it is not necessary to use formulas intended to adulate, for base propagandistic and electoral ends, the spirit of sedition of the strata inferior to the masses disguised as ``social justice". In general, the same style of active impersonality, dignity, solidarity in the production that is typical to the ancient professional and artisan corporations should be recovered. Syndicalism, with its ``battle" and those genuine threats which it now offers us too many examples, is to be banished. But, let us repeat, much one must reach that point starting from the interior. The important thing is that against every from of resentment and social antagonism everyone knows how to recognize and love his own place, that conforms to his own nature, recognizing thus even the limits between which he can develop his possibilities and follow his own perfection: because an artisan who perfectly discharges his function is undoubtedly superior to a king who is unfit and not at the height of his dignity.

In particular, we can permit a system of expertise and corporative representation, to supplant the parliamentarianism of the parties; but we should keep in mind that the technical hierarchy, in their entirety, can signify nothing more than a level in the integral hierarchy: they concern the order of means, to be subordinated to the order of ends, to which only the properly political and spiritual part of the State corresponds. To speak instead of a ``State of workers" or of production is to make the part equivalent to the whole, the same as reducing the human organism to its simply physical and vital functions. Nor can a similar obtuse and dark thing be our emblem, nor the ``social" idea itself. The true antithesis facing both the ``East" and the ``West" is not the ``social ideal". It is instead the \emph{integral hierarchal idea}. In respect to that, nothing uncertain is acceptable.

\paragraph{Point 7}
\begin{quotex}
Highlights of point 7:

\begin{itemize}
\item
The ideal is a virile and organic political unity 

\item
This unity is fundamentally a spiritual unity 

\item
The dignity and freedom of the human person is found in an organic society, not in individualized liberalism 

\item
Evola rejected the Salo Republic because it excluded the Italian monarchy 

\item
The spiritual environment must be created prior to the political and a suitable symbol created 

\end{itemize}

The Carlyle quote is from \textit{Heroes and Hero Worship} and is misquoted in \textit{Men Among the Ruins}. This is the full quote:

The Valet-World has to be governed by the Sham-Hero, by the King merely \emph{dressed} in King-gear. 

\end{quotex}
If the ideal of a virile and organic political unity was an essential part in the world that had to be swept away—and through it, we also brought back the Roman symbol—we must then recognize the cases in which such a need deviated and was almost miscarried in the mistaken direction of \emph{totalitarianism}. This, again, is a point that must be seen with clarity, so that the differentiation of the political fronts is precise and, also, arms are not furnished to those who want to confuse things after due consideration. Hierarchy is not hierarchism (an evil that, unfortunately, today at times tries to come back in minor tones), and the organic conception has nothing to do with a statolatrous sclerosis and a leveling centralization. As for the individual, a true surpassing both of individualism and collectivism happens only when men are face to face with men, in the natural diversity of their being and their dignities. And as for the unity that must prevent, in general, every form of dissociation and absolutization of the particular, it must be essentially spiritual, it must be a central, orienting influence, an impulse that, depending on the leaders, assumes very differentiated forms of expression. This is the true essence of the ``organic" conception, opposed to the rigid and extrinsic relations typical of ``totalitarianism". In these fields the requirement of dignity and liberty of the human person that liberalism can conceive only in individualistic, egalitarian, and privatistic terms can be realized integrally. It is in this spirit that the structure of a new political-social ordering must be studied, in solid and clear articulation.

But such structures need a center, a supreme point of reference. A new symbol of sovereignty and authority is necessary. The delivery, in this regard, must be precise and ideological prevarication cannot be admitted. It is good to say clearly that here we are only talking about the subject of the so-called institutional problem; it is first of all about what is necessary for a specific \emph{climate}, for the fluid that must animate every relationship of fidelity, dedication, service, disindividualized action, so that the gray, the mechanistic, and the oblique of the current political-social world is truly surpassed. Here today it will therefore end up with no way out when at the top it is not capable of a type of ascesis of the pure idea.

Both the less felicitous antecedents of our national traditions, and even more, the tragic contingencies of yesterday's prejudice are, for many, the clear perception of the right direction. We can still recognize the inconclusiveness of the monarchical solution, when those who today can only defend a residue of the idea are kept in view, a symbol emptied and devitalized, which is that of the parliamentary constitutional monarchy. But in an altogether decisive way the incompatibility in relation to the republican idea must be pointed out. To be anti-democratic on the one hand, and on the other to defend ``ferociously" (this is unfortunately the terminology of some exponents of a false intransigence) the republican idea is an absurdity that is the proof of it: the republic (we mean modern republics: the ancient republics were of the aristocracy, as in Rome, the latter of the oligarchy, often with the character of tyranny) belong essentially to the world that rose to life through Jacobinism and the anti-traditional and anti-hierarchical subversion of the XIX century. And to such a world, which is not ours, let it go. In terms of principle, a currently monarchical nation that becomes a republic can only be considered as a ``degraded" nation. For Italy one does not play in the ambiguity in the name of a fidelity to Fascism of Salo, because if, for this reason, one should follow the false republican way, at the same time it would be unfaithful to something more and better, it would throw into the sea the central nucleus of the ideology of the 20s, i.e., its doctrine of the State in regards to authority, power, \emph{imperium}.

It is necessary to conserve this doctrine, only without agreeing to level down and without making it the game of any group. The concretization of the symbol can be left indeterminate for now; the essential task is to prepare silently the spiritual environment adapted to what the symbol of an intangible elevated authority has felt and reacquired the fullness of its meaning: it cannot correspond to the stature of any revocable ``president" of the republic, and not even that of a tribune or populist leader, the possessor of mere individual formless power, lacking every higher blessing, supporting himself instead on the precarious prestige he exercised on the irrational forces of the masses. It is that to which some have given the name of ``Bonapartism" and that was justly recognized in its meaning not as the antithesis to demagogic or ``popular" democracy, but also as its logical conclusion: one of the dark apparitions in the Spenglerian ``decline of the West". There is a new touchstone for us: the \emph{sensibility} in respect to all that. Carlyle had already written: ``the Valet-World has to be governed by the Sham-Hero" and not by a real Ruler.

\paragraph{Point 8}
\begin{quotex}
In this point, \textbf{Julius Evola} brings up the topic of nationalism, which he rejects. In particular, there is no support for anything calling itself ``White Nationalism", despite those movements that claim to be ``radical traditionalist" or ``influenced" by Evola. That would be to be based on ``naturalistic" principles, a notion totally opposed to any Traditional state.

Rather, it must be based on the ``Idea", and of course, this Idea is the Traditional idea. This is not at all comparable to the American claim to be a ``proposition nation", that is, it is based on an ``idea", not on an ethnicity, religion, or culture. The American proposition is that ``all men are created equal", the polar opposite of the hierarchical state. In theory, that means that each man decides what is true and good, and determines his own meaning of life. However, what that means in practice is the idolization of the least common denominator. In theory, each man would be free to actualize his highest possibilities; in practice, the harsh reality is that the possibilities of the mass of men are quite limited. In the Traditional State, dominated by the highest spiritual ideal, common men are raised beyond their natural capacities and find their true meaning in higher values.

We see Evola proposing an Order, presumably along the lines of the Templars, that is, an order that is both spiritual and chivalric. That would only be possible when there are a sufficient number of men who both know and are faithful to higher principles, and can testify as witnesses to a higher authority. 

\end{quotex}
Another point must be made in an analogous order of ideas. It concerns the position to take in regards to nationalism and the general idea of homeland. That is so much more timely, inasmuch as today, in order to try to save whatever possible, many people would like to adopt a sentimental and, at the same time, naturalistic conception of the nation, notions foreign to the highest European political traditions and irreconcilable with the same idea of the State which we spoke about. Even prescinding from the fact that we see the idea of homeland being invoked rhetorically and hypocritically by the most opposed parties, even by exponents of red subversion, that conception is not now factually at the height of the times because from one side it assists the self-formation of great supranational blocks, and from the other it always appears more necessary to find a European point of reference, unifying beyond the inevitable particularism that inheres in the naturalistic idea of the nation and still more in ``nationalism".

Nevertheless, the question of principle is more essential. The political plane in such regards is that of an elevated unity in respect to the unities defined in naturalistic terms as are also those to which the generic notions of nation, homeland, and people correspond. In this higher plane, the idea unifies and divides, an idea carried by a resolute elite and tending to be concretized in the State. For this reason the Fascist doctrine—that part of it which remains faithful to the best European political tradition—gave primacy to the Idea and State in respect to nation and people and it understood that nations and people acquire a meaning and form and participate in the higher degree of existence only within the State. Even in periods of crisis, as there is currently, it is necessary to hold firm to this doctrine. Our true homeland must be recognized in the Idea. Not being from the same land or the same language, but being of the same idea is what counts today. This is the base, the point of departure. As for the collectivistic unity of nation, of the children of the homeland\footnote{In French, \emph{des enfants de la patrie} from \emph{La Marseillaise}.}—which has predominated more and more from the Jacobin revolution until now, we in any case oppose it with something like an \emph{Order}, men faithful to principles, witnesses to a higher authority, and legitimacy proceeding precisely from the Idea.

In the matter of practical ends today it is desirable to arrive at a new national solidarity, only we do not descend to compromises in order to reach it; the presupposition, without which every result would be illusory, is separation from that and taking the form of a definite political arrangement from the Idea—such as political ideas and visions of life. Even today there is not another way: it is necessary that among the ruins the process of the origins be renewed, that which, depending on elites and a symbol of sovereignty or authority, unites the people in the great traditional States, as forms arising out of the unformed. This realism of the idea is not to be understood as meaning to remain on a fundamentally sub-political plane: that of naturalism and sentimentalism, if not completely of the patriotic rhetoric.

And in case we want to support our idea also on national traditions, one remains very careful: because there is an entire ``homeland history" of Masonic and anti-traditional inspiration specializing in attributing national Italian character to the most problematic aspects of our history: starting from the revolt of the Communes supported by Guelphism. With it, a tendentious ``Italianity" takes prominence in which we cannot and do not want to recognize ourselves. We leave it voluntarily to those Italians who celebrated the ``second Risorgimento" with ``liberation" and partisanship.

Idea, order, elite, State, men of Order—in such terms the lines are maintained, as long as it be possible.

\paragraph{Point 9}
\begin{quotex}
In Point 9, \textbf{Julius Evola} rejects the ``solutions" offered by the modern world. First of all, it is not simply a matter of preserving a certain ``culture", which is the superficial, but rather of developing an entire ``worldview", in the heart of man's being, from which culture will arise as an effect.

This worldview does not come from reading books and, we should add in our age of mass media, it definitely cannot arise from passively absorbing the ideals, values, and opinions of main sources of public discourse today. The ``right" worldview cannot only arise within, from living a noble life, perhaps following, as Evola said in different circumstances, the \emph{Path of the Gods}. Personally, I find that the Noble Eightfold Path of the Buddha expresses this idea: right view, right aspiration, right speech, right action, right livelihood, right effort, right mindfulness, right concentration.

Evola then singles out three modern worldviews for a withering critique: darwinism, pschoanalysis, and existentialism. All three—and even today they form the foundation for the modern and postmodern wordlview albeit in somewhat mutated forms—reduce man to the lowest level. He is an ``evolved" animal, he is the plaything of unconscious forces, he is dominated by feelings of angst and meaninglessness. Anyone who opposes these worldviews is then branded as ignorant or even evil. As he points out, the worst invective is directed to those who refuse to accept darwinism as a valid understanding of the origin of man; the darwinian dogma is enforced more strictly than any previous Inquisition. 

\end{quotex}
Something needs to be said about the problem of culture. Not beyond measure. We in fact do not overvalue culture. What we call a ``worldview" is not based on books; it is an inner form that can be more precise in one person without a particular culture than in an ``intellectual" or a writer. We must ascribe the fact that the individual is left open to influences of every type to the weight of everything among the bad omens of ``free culture", even when he is not able to be active in relation to them, to know how to discriminate, and to judge according the right judgment.

Thus discourse cannot just point out that, as things currently stand, there are specific currents from which the youth of today must defend themselves interiorly. We spoke initially of a style of rectitude, of inner resistance. This style implies a right knowing and young men especially must take account of the intoxication working on the totality of a generation from concordant varieties of a distorted and false vision of life, that have affected men's internal strength. In one form or another, these toxic elements continue to act in the culture, in science, in sociology, in literature, as so many points of infection that must be identified and beaten. Aside from historical materialism, which we already mentioned, Darwinism, psychoanalysis, and existentialism are among the major ones.

The fundamental dignity of the human person must be asserted against Darwinism, recognizing his true place, which is not that of a particular, more or less evolved, type of animal among so many others, differentiated through ``natural selection" and always tied to bestial and primitivistic origins, but he is to elevate himself virtually beyond the biological plane. If today no one speaks much anymore about Darwinism, its substance nevertheless remains. The biological Darwinian myth, in one or another of its variants, has the precise value of a dogma, defended by the anathemas of ``science", in materialism, whether of the Marxist or American civilization. Modern man is addicted to this debased idea, he recognizes himself in it comfortably and finds it natural.

Against psychoanalysis the ideal of a Self (or ``I") must be validated, a Self who does not abdicate, who intends to remain conscious, autonomous, and sovereign in the confrontation with the dark and subterranean part of his soul and the demon of sensuality; that does not mean either ``repression" or a psychotic scission, but he brings about an equilibrium of all his faculties ordered to a higher meaning of living and acting. An obvious convergence can be pointed out: the disempowerment of the conscious principle of the person, the prominence given to the subconscious, the irrational, the ``collective unconscious" and similar ideas from psychoanalysis and analogous schools, corresponds in the individual precisely to what the urgency, the impulse from below, the subversion, the revolutionary substitution of the inferior to the superior, and the disdain for every principle of authority, represent in the modern social and historical world. The same tendency acts on two different planes, and the two effects cannot avoid being mutually integrated.

As for existentialism, even to distinguish in it what is properly a philosophy—a confused philosophy—it remained, until our time, of pertinence to restricted circles of specialists. It is necessary to recognize the state of soul in crisis that became a system and was adulated, the truth of a broken and contradictory human type who experiences, with anxiety, tragedy, and the absurd, a freedom from which he does not feel elevated, and to which he instead feels without escape and responsibility, condemned in the midst of a world deprived of value and meaning. All this, when the best Nietzsche had already indicated a way to recover a meaning for existence and to give to oneself a law and an intangible value even in the face of a radical nihilism, in the sign of a positive existentialism, according to his expression: from a ``noble nature".

Such are the lines of surpassing that must not be merely intellectualized but lived, realized in their direct meaning for the interior life and one's own conduct. To lift oneself up is not possible as long as one remains as though he is under the influence of similar forms of false and deviant thinking. Detoxified, a man can pursue clarity, rectitude, strength.

\paragraph{Point 10}
\begin{quotex}
Julius Evola points out the three possible reactions to the modern world, which he regards as fundamentally ``bourgeois". We may dispute that today in that proletarian values seem to predominate in our time.

\begin{enumerate}
\item
Oppose the bourgeoisie with a ``collectivized and materialized humanity," to which we may add an ``undifferentiated" humanity. 

\item
Maintain the bourgeois status quo through focusing on the economic life, solely. 
\item
Oppose the bourgeoisie with a ``higher heroic and aristocratic conception of existence", which is both ``anterior and superior" to the current state of affairs. 

\end{enumerate}
\end{quotex}
The position that lies in the zone between culture and custom will finally be made clear. From communism came the rallying cry of the anti-bourgeois that was also picked up in the cultural groups of certain ``committed" intellectual circles. As bourgeois society is something intermediate, so there exists a double possibility of surpassing the bourgeois, of saying ``no" to the bourgeois type, to bourgeois civilization, to the bourgeois spirit and bourgeois values. One corresponds to the direction that leads even lower than all that, toward a collectivized and materialized humanity with its Marxist ``realism": social and proletarian values against bourgeois and capitalist decadence. But the other is the direction of those who oppose the bourgeoisie in order to actually rise up beyond it. The men of the new arrangement will still be anti-bourgeois, but by way of the aforesaid higher heroic and aristocratic conception of existence; they will be anti-bourgeois because they disdain the comfortable life; anti-bourgeois because they will follow not those who promise material advantages, but those who demand everything from themselves; anti-bourgeois, finally, because they are not preoccupied with security but love the essential union between life and risk, on all planes, making the inexorability of the bare idea and precise action his own. Yet another aspect, through which the new man, cellular substance for the motion of awakening, will be anti-bourgeois and will differentiate himself from the preceding generation, is through his intolerance of every form of rhetoric and false idealism, of all those great words that are written in capital letters, of everyone who is merely a gesture, words intended for effect, scenography. Essentiality, instead, a new realism in pitting itself exactly against the problems that will be imposed, in doing so the value is not appearing, but rather being, not prattling, but rather realizing, in a silent and exact way, in sync with similar forces and in adherence to the command that comes from above.

Those who are able to react against the forces of the left only in the name of idols, of lifestyle and the mediocre conformist morality of the bourgeois world, have already lost the battle beforehand. This is not the case for the man who stands on his feet, having already passed through the purifying fire of internal and external destruction. This man, in the same way that politically he is not the instrument of a false bourgeois reaction, so, in general, he regains forces and ideals anterior and superior to the bourgeois world and the economic era, and it is with them that he creates the lines of defense and consolidates positions when, at the opportune moment, he will strike out with the action of reconstruction.

Even in such regard, we intend to reclaim a delivery not followed: because we know how in the Fascist era there was an anti-bourgeois tendency that had tried to explicate itself in a not dissimilar sense. Unfortunately even there the human element was not up to the task. And they knew how to make rhetoric even out of anti-rhetoric.

\paragraph{Point 11}
\begin{quotex}
In the final point, \textbf{Julius Evola} addresses the question of spirituality. He first considers Catholicism, the last Tradition in the West, but rejects it as inadequate in our time because the Church no longer represents that Tradition. If, instead, had it maintained that Tradition, had it still represented the only force against the modern world, and had it held firm to the principles of \textit{The Syllabus Of Errors}\footnote{\url{https://www.papalencyclicals.net/Pius09/p9syll.htm}} (which everyone should read), Evola would support it.

Evola does not propose some other tradition, alien to Europe, but rather counts on new men, aware of the transcendent, aware of life beyond death, to forge a new spirituality, in expectation of a blessing from above. Given his anti-clerical attitude, this is a somewhat inconsistent position. What is there to prevent such men from recovering that tradition on their own? Why must they wait passively for the ``church" to give it to them? The Kingdom is taken by force. 

\end{quotex}
Let us briefly consider a final point, that of the relationship with the dominant religion. For us, the secular State, in whatever form, belongs to the past. And, in particular, we oppose its parody that asserted itself, in certain circles, as the ``ethical State", the product of a weak, spurious, empty ``idealistic" philosophy associated with Fascism [that of Giovanni Gentile], but through its nature such to give equal approval, like a dialectical game of dice, to the anti-Fascism of a [Benedetto] Croce.

Yet though we oppose similar ideologies and the secular State, a clerical or clericalistic State is likewise unacceptable for us. A religious factor is necessary as the background for a true heroic conception of life that must be essential for our political alignment. It is necessary to feel in oneself the evidence that there is a higher life beyond this terrestrial life, because only those who feel this way possess an unbreakable and unconquerable strength, only they will be capable of an absolute enthusiasm. But where this is lacking, to defy death and to take no account of one's own life is possible only in sporadic moments of elation or in the unleashing of irrational forces: nor is there a discipline that can be justified, in the individual, with a higher and independent meaning.

But this spirituality, which must be alive among us, does not require obligatory dogmatic formulations of a given religious confession; however, the style of life that must be drawn from it is not that of Catholic moralism, which aims for little more than a virtuistic domestication of a human animal; politically, this spirituality cannot but foster mistrust in respect to everything that as humanitarianism, equality, principle of love and of the forgiveness instead of honor and justice; it is an integral part of the Christian conception. Certainly, if Catholicism were capable of making a program of high ascesis its own and exactly on this base, almost like a recovery of the spirit of the best Medieval crusader, makes of faith the soul of an armed bloc of forces, almost a new close-knit Templar Order, relentless against the currents of chaos, breakdowns, subversion, and the practical materialism of the modern world—certainly, in such a case, and also in the case that, at the minimum, it held firm to the positions of the Syllabus, there could not be a single instance of doubt about our choice.

However, as things currently stand, given the mediocre and, fundamentally, bourgeois and parochial level, to which today practically everything that is confessional religion has descended and given the modernist destruction and the rising opening to the left of the post-conciliar Church of the ``aggiornamento", for our men, pure reference to the spirit will suffice, precisely as the evidence of a transcendent reality, to appeal to another force in order to graft onto our force, to have a foreboding that our battle is not only a political battle, to attract an invisible consecration on a new world of men and leaders of men.

\paragraph{Conclusion}
\begin{quotex}
With this paragraph, \textbf{Julius Evola} concludes his Orientations. Note that he asks the next generation to take up the torch, but not necessarily to repeat exactly the past. Lessons must be learned and Tradition expressed perhaps in new terms. Fighting old battles may not be helpful in the present battles.

Note that Evola, the ``man of action", insists on intellectual rigour and intransigence with respect to the ideas of Tradition. Those who hold firm to those ideas will be united. This unity cannot be achieved on the basis of some contingent political position; certainly not with those who deviate from such ideals while claiming to be ``traditionalists". Those who keep their minds focused on the ideals of Tradition will never be overcome by the enemy. They remain firm in the conviction that ``what can be done, will be done."

As a side note, please note the reference to ``illusory order". This distinction will be important in the discussion of the Logos. 

\end{quotex}
These are some essential orientations for the battle to be fought, especially in respect to the youth, who take up the torch and the baton from those who have not fallen, learning from the errors of the past, knowing how to discriminate and review everything that was affected, and still today is affected, by contingent situations. It is essential not to descend to the level of the adversaries, not to be reduced to repeating the same mantras, not to insist unduly on those of yesterday which, even if worthy of being remembered, do not have the impersonal and current value of an idea-force, not to yield to the suggestions of false politicizing realism, the mark of every ``party". It is necessary for our forces to act also in political hand-to-hand combat in order to create all the possible space in the current situation, and to contain the assault, otherwise almost uncontested, of the forces of the left. But beyond that it is important, it is essential that an elite constitute itself which, in a gathering intensity, identified by an intellectual rigour and an absolute intransigence, depending on the idea that must be united and affirmed above all in the form of the new man, of the man of the resistance, of the man standing up among the ruins. If the task will be to go beyond this period of crisis and vacillating and illusory order, the future will be up to this man. But when also the destiny — that the modern world created for itself and is now sweeping away — should not be restrained, beside such remarks the interior positions will be preserved: in whatever eventuality, what can be done will be done and we will belong to that country which will never be occupied or destroyed by any enemy.


\hfill



\flrightit{Posted on 2012-07-01 by Aeneas }

\begin{center}* * *\end{center}

\begin{footnotesize}\begin{sffamily}



\texttt{apeiron on 2012-07-10 at 19:06 said: }

``Starting from what can still remain among the ruins, to reconstruct slowly a new man to animate through a determinate spirit and a suitable vision of life, to fortify through the tenacious adherence to given principles – this is the true problem."

This problem is a very true one that still begs for a solution and surely it will be ongoing for the foreseeable future. For it appears that each rectitude since the end of WW2 has been one of a personal nature rather than one holistically aligned. I praise Gornahoor for trying at least to bridge this divide in an original way by continuing Dantes' task, the hermeticism and templar Rosicrucianism of transcendent European sovereignty that links a Tradition from our deepest hyperborean antiquity to our break with the present circumstances. Guenons example to abandon his own Tradition and seek that of an Islamic one, is not ours and neither can Evolas' personal equation amount to a solid example for emulation despite the greatness of the man's works. Yes I'm generalizing in this post without getting into academic specifics, but its holds true all the same.

To quote Mouravieff, from Gnosis (book III):

``Since then, the revolutions that have successively brought us steam, electricity, mechanimtion and atomic energy have transformed the face of the world, and habit as a substitute for buppiness has disappeared during these upheavals. This means that, although we hear people claim happiness with growing fewour in the feverish atmosphere of our times, we find ourselves in a void. If a few `last of the Mohicans,kith their outdated mentality, raise their voices to call for a return to the `normal' state of things, in their ndive sincerity they appear like mediaeva1 knights leading a cavalry charge against armoured cars" 

It is slightly puzzling to understand why such an important work as ``Orientamenti" had not been translated into English before so for that reason I extend my thanks to Gornahoor for taking the initiative. If all systems have fallen and all ideologies have failed then surely a guide of this sort that works on interior strengthening of the spirit is needed more than ever. At the same time I'm not suggesting anything conclusive will be drawn from this, but I look forward to the subsequent points in this series. We know what the problem is but the answers are lacking as usual, and even the best minds spend years in the wilderness piecing together the pieces of a fragmented interior.


\hfill

\texttt{Dominion on 2012-07-15 at 12:29 said: }

Seeing at last one of the few Traditionalist writings in which we get a view of how the revolt takes place is definitely inspiring. Evola's idea that the Traditional man must simply live as a Traditional man, and make himself known by virtue of his existence is simple and yet highly practical. Assuming that a personal inner order has been achieved, or is being achieved, it lets this fact be transmitted into the real world. It can be done in daily life, at any level of society. In a sense, the man of Tradition then displaces the current order by refusing to be a part of it through his personal actions. It reminds one of Michael Collins' dictum that the Irish would ``defeat the British Empire by ignoring it." Thanks for publishing.


\hfill

\texttt{Logres on 2012-08-05 at 20:59 said: }

Baron Ledhin proposed something along the lines of corporativism:

\url{http://www.phillysoc.org/Portland.htm}


\hfill

\texttt{Cologero on 2012-08-06 at 07:18 said: }

Logres, that is one ugly hard-to-read web site.

To consider the question about ``why" to deal with such political/economic topics, particularly since it can only be done in a Platonic way with no chance of implementation. First is as an antidote to the many web sites that claim to be ``radical traditionalists", ``influenced" by Evola, when they clearly have no such program as that envisioned by Evola, often run by people who ``no longer know what normal is".

The hierarchy of spiritual values, then military/political, and lastly economic is the traditional social arrangement, described even by Hermetic authors. Can such an arrangement even be considered ``in thought" at this time? Clearly, this arrangement has been the target of overthrow by all revolutionary movements. To this day, the ideal of religious authority, martial spirit, and cooperative economic activity is falsely considered ``fascist", and every vestige of it is attacked. For example, an attack on the remnants of spiritual authority is disguised as a women's health issue. The warrior spirit is mitigated by egalitarian pressures in the armed services. And so on … why would there not be any value in considering such events from the standpoint of Tradition?


\hfill

\texttt{Logres on 2012-08-06 at 09:51 said: }

``The important thing is that against every from of resentment and social antagonism everyone knows how to recognize and love his own place, that conforms to his own nature, recognizing thus even the limits between which he can develop his possibilities and follow his own perfection: because an artisan who perfectly discharges his function is undoubtedly superior to a king who is unfit and not at the height of his dignity."

Great quote. Ledhin used to say ``loving command, and proud obedience" were the watchwords of the organic Catholic order.


\hfill

\texttt{logres on 2012-08-12 at 14:57 said: }

I must say that Evola always surprises me with what a precise mind he possessed – once he says it, very often, it clears away a century of thought, and years of personal confusion:

``In these fields the requirement of dignity and liberty of the human person that liberalism can conceive only in individualistic, egalitarian, and privatistic terms can be realized integrally.."

This was exactly the same goal aimed at by Baron Ledhin, who (unfortunately) did not understand the abiding ideal of imperium – once the monarchs were all equal, the World Wars commenced, and they seemed powerless to stop it (though they personally wished to, at some level). I wonder if Ledhin knew of Evola? I'll have to check…


\hfill

\texttt{logres on 2012-08-12 at 15:08 said: }

" Indeed, Kuehnelt-Leddihn observed that one of the weaknesses of the Right was its failure to failure to offer a utopia of its own as a counter to the utopias of the Left.9[2033?] (Leftism Revisited)

Indeed, ``Erik von Kuehnelt-Leddihn was fond of pointing to the traditional Chinese civil service examination system as one reflecting the ideals of meritocracy. The scholars who comprised the Mandarin class were drawn from the ranks of those demonstrating superior skill and ability. This was not a hereditary class but one where even the lowliest peasants with remarkable talents could achieve self-advancement."

Interesting that Ledhin (like CS Lewis) seems to have appreciated the possible directions ``things might have gone", but appeared powerless to actually articulate it. One wonders what might have happened had a more coordinated, concerted, and spirited effort been made to re-constitute the opposite of a Revolution, without falling into the Benoist (or similar) side-tracks.


\hfill

\texttt{Cologero on 2012-08-12 at 20:08 said: }

Good points, Logres. I was speaking to someone a couple of weeks ago, who claims to follow the New Right. I asked about his (or their) vision. He envisions communities of pagans, practicing their own way, apparently like the city states of Ancient Greece. I asked him the a question, equivalent to what you mentioned about kings. How do you guarantee the ``kings" won't begin to battle each other and consolidate; especially since they want to be a ``virile" (at least some do) and warrior society. I also pointed out that they would have no way to protect themselves against the dreaded ``monotheists", who would probably have the imperium in mind. Curiously, he tried to shame the monotheists for being intolerant; sounded just like a liberal.

As Evola points out, the spiritual vision comes first, the symbol to gather around. It will be a hard sell, since liberal individualism is so ingrained and anything else seems a threat to that willfulness that is called ``freedom". One of the dominant spiritual visions involves being raised up into the clouds; doesn't make sense and is untraditional, yet it holds a lot of hearts and minds.

A meritocracy like the Mandarins is hardly possible now, as we know who will administer the tests, what answers will be deemed acceptable, and who will get ``bonus" points. I suppose we are already, in a way, ruled by Mandarins, just not what Kuehnelt-Leddihn would have hoped for.


\hfill

\texttt{apeiron on 2012-08-12 at 20:35 said: }

I always resonated with the simplicity inherent in the spiritual unity of the true golden Imperium, everything else just falls into place almost with a non-rational(supra-rational) gnosical understanding. 

``One of the dominant spiritual visions involves being raised up into the clouds; doesn't make sense and is untraditional, yet it holds a lot of hearts and minds."

Which dominant spiritual vision do you speak of here? Christianity, Judaism or Islam? Are these untraditional visions in your eyes?


\hfill

\texttt{Cologero on 2012-08-12 at 21:57 said: }

Apeiron, apparently you are not familiar with religious trends in the USA. I was referring to the neo-Christians in North America who believe they will be lifted up into the clouds before the end of the world as described here, for example\footnote{\url{http://en.wikipedia.org/wiki/Dispensationalism}}. There are many such believers in the USA. Not every Christian sect is ``traditional" in its outlook.

The point is that there is no spiritual vision that can unify the West at the present time. Part of task pointed out in Point 7 is to see clearly the various movements competing for that vision.


\hfill

\texttt{Dominion on 2012-08-12 at 21:59 said: }

A similar system still exists in China with the CPC's cadre apparatus. Prussia's Junker aristocracy is another good example of a sort of ``aristocracy of merit".

This statement by Evola is, I think, very handy both as a repudiation of totalitarianism and aspects of historical fascism, and a guide to what exactly was wrong with it. When we look at fascist movements, there is a very big pattern of trying to preserve traditional forms without adhering to the truths that the forms are meant to embody. For example, Italian fascism's focus on the Roman Empire in its expansionist policy and ``Caesarism", rather than in its spiritual Tradition. Spanish fascism's focus on the Catholic Church and the archetypal peasant, which resulted in the stagnation of many aspects of Catholic life in that country as the State used the Church as a justification for its totalitarian rule (writer Richard Wright reports that fornication among the Spanish agricultural classes was higher than in most other ``liberal" European countries, a pattern that continues today in the American South with things like teenage pregnancy and use of pornography). National Socialism's focus on the racial Aryan, and not the spiritual one. 

It is of utmost importance to realize the steps in which a transition and reformation in ideas must go. The ideas by some groups to implement an ``Imperium" in Europe or the West at this stage is nonsensical. The first great task is to enable individuals to grasp the truths of Tradition themselves and begin their `personal' transformation. As these numbers grow, it may become possible to create bigger and bigger networks. Then, and only then, can one even begin the think about restoring Tradition at the level of society and State. It may be well for us to purge the idea of ever seeing these later stages in our lifetimes, for it would be better to die knowing that a rebirth is being carried out step by small step than to see a false Imperium created by those who cannot see that it is a form containing nothing.


\hfill

\texttt{Ferdinand Gombault on 2012-08-13 at 09:59 said: }

logres, Erik von Kuehnelt-Leddihn knew of Evola; in one of his works he mentions his name in a somewhat condescending manner pointing to his ``paganism". If you wish, I could try to find the exact quotation. However, I doubt that Kuehnelt-Leddihn was really familiar with Evolas writings.

The two had a decisive dissent: Kuehnelt-Leddihn was a fervent catholic and I do not have to explain Evolas conviction in this regard. On top of that, Kuehnelt-Leddihn regarded himself as a ``liberal" (though in a Central-European sense, which does not have the meaning ``left", for Kuehnelt-Leddihn rather to the contrary), which I strongly doubt Evola would have done. Kuehnelt-Leddihn also was decidedly anti-rascist and anti-fascist, although anti-egalitarian and ``right-wing". I personally would recommend reading his books as a stimulating contrast to Evola.


\hfill

\texttt{Logres on 2012-08-19 at 16:31 said: }

Wouldn't there be a difference between arguing for a ``white" Imperium \& prudentially believing that certain ethno-states be maintained in relative genetic stability, provided that the entity was not allowed to exalt itself to a transcendental level, or expand at the expense of other, similar entities? Or is that an inevitable betrayal of principle, as well? Is the distinction moot at this point?


\hfill

\texttt{apeiron on 2012-08-19 at 19:15 said: }

It is fair to say that Evola rejected the Jacobinic manifestation of the nationalist idea as lesser to that of a spiritual Imperium of Aryan races, hierarchically aligned. Why National Socialism and Fascism were not explicitly called ``white" or ``Aryan" nationalism is simply because there was no demographic crisis of white peoples in their indigenous homelands or of their genetic stock worldwide: the first tier of race. At least back in the thirties one had, more or less, only the spiritual problem to worry about.

In Evolas words, from ``Doctrine of Awakening": 

``We have to remember that behind the various caprices of modern historical theories, and as a more profound and primordial reality,

there stands the unity of blood and spirit of the white races who created the greatest civilizations both of the East and West, the Iranian and Hindu as well as the ancient Greek and Roman and the Germanic."

To say that there was no support for ``anything calling itself ``White Nationalism"" may be true but slightly misleading. Equally one could make the same argument for National socialism and Fascism, both modern themes influenced by masonic ideas yet stood as the last bridge-stone since the middle ages to surpass modernism with the right direction. Both of these regimes Evola worked with, had they not been crushed by two daemonic hyper-material superpowers, it is only our speculation what may have come of them.


\hfill

\texttt{Cologero on 2012-08-19 at 20:39 said: }

Logres, Evola's context was post-WWII Europe; he was pointing a way out of the ruins. Hence, a ``white" imperium was implicit, given the ethnic constitution of Europe in 1951; he was not speaking of any other situation or other civilizations. The point, nevertheless, is that race, or even ethnicity, is insufficient as a unifying principle, and the facts bear that out. By definition, an Empire consists of ethno-states. So what Idea would unite them? Certainly, ethno-nationalism in itself would not.

As for Apeiron's thoughtful point, we wonder if the race that created the ``Iranian and Hindu as well as the ancient Greek and Roman" civilizations even still exists; certainly that spirit no longer exists. Ironically, this is implied in his comment about the ``two daemonic hyper-material superpowers", the modern equivalent of the ``civilizations both of the East and West".


\hfill

\texttt{Sparrow on 2015-03-02 at 03:01 said: }

I once made the point that White Nationalists with their petty agendas, orgiastic violence, and inability to contemplate the transcendent, had more in common both psychologically and in spirit with a drunken Red Army officer than the calm, disciplined SS foot soldier they desire to emulate. The remark was not well received, but I feel it had a certain truth to it.


\hfill

\texttt{Michael on 2012-09-02 at 20:46 said: }

Will the action of reconstruction happen spontaneously, or will it be the result of an organization? In other words, is it just a matter of living in a certain way (riding the tiger) or the result of a movement?


\hfill

\texttt{Cologero on 2012-09-04 at 12:09 said: }

Michael, nothing happens ``spontaneously".

First, the conscious and unconscious agents of revolution need to be identified

Then, of the former, the covert and open agents of revolution

Then there must be an understanding of what revolution is

Then there will be the understanding of the opposite of the revolution

The possible is the real, hence the conditions to make the foregoing possible must be identified.


\hfill

\texttt{Cologero on 2012-09-04 at 20:23 said: }

There is a rather odd debate\footnote{\url{http://theinteramerican.org/blogs/olavo-de-carvalho/257-olavo-de-carvalho-debates-alexandr-dugin-ii.html}} between Olavo de Carvalho and Alexander Dugin. In Section 9 Carvalho describes the possible historical agents. Although both men are Traditional men of the Right, influenced by Guenon, they are far apart from each other. I don't believe either one identified fully the revolution; whether their solutions are fully convincing is a matter for debate.


\hfill

\texttt{Logres on 2012-09-05 at 05:53 said: }

That was an odd debate – I found myself in sympathy with the Westerner, but nodding at points Dugin was making, such as the necessity to love both a country/culture's flaws as well as excellencies (in a sense). Any comments you care to make further on this debate would be much appreciated, Cologero. Dugin has a Gnostic streak, at least in rhetoric, that is difficult to contextualize, for me.


\hfill

\texttt{Jason-Adam on 2012-09-05 at 06:17 said: }

Not to go tangental, but I think that in order to fully understand what Olavo and Dugin were debating, we first need to consider the ideas of Franco Freda, the follower of Baron Evola, and his book THE DISINTEGRATION OF THE SYSTEM. I am sure Cologero has read the book, but to summarise it, Freda's argument was that men of Tradition are loyal to timeless principles irrespective of race or culture and thus if the West has become anti-Traditional beyond all hope of recovery, as Evola argued in Ride the Tiger, to restore a Traditional world it is first necessary to destroy the West; hence Freda's call for the ``Right" to ally with the extreme and violent left. This is a huge influence on Dugin's worldview.


\hfill

\texttt{Michael on 2012-09-05 at 20:10 said: }

When you said agents of revolution, I thought of the individuals and organizations that initiated and spread Enlightenment ideas, but Carvalho appears to have a broader understanding of what an agent is.

If Carvalho is correct, then the only hope of a return to Tradition in the West is the Catholic Church since it is doubtful that the other entities named (royalty \& nobility, initiatory \& esoteric organizations) are willing or able to do anything. This is especially true if Guenon was right in saying that initiatory organizations no longer exist in the West.


\hfill

\texttt{Cologero on 2012-09-05 at 23:52 said: }

On one point, the two men agree: there is a globalist financial elite. Carvalho supports a conspiracy theory, identifying that group as the ``Syndicate", transnationalist, not confined to any particular country. Dugin, OTOH, conflates that group with American power \emph{in toto}. Whence, their widely divergent views. Carvalho\footnote{\url{http://theinteramerican.org/blogs/olavo-de-carvalho/257-olavo-de-carvalho-debates-alexandr-dugin-i.html}} identifies three movements that can oppose the Syndicate: (Section 5)

\begin{enumerate}
\item An international union of Catholic and Protestants 
\item Zionism 
\item American neo-conservatism 
\end{enumerate}
There is no evidence of them acting in concert, nor even verbally opposing the plans of the ``Syndicate", not even in muted tones. Nevertheless, Carvalho wishes to defend the existence of Catholic-Protestant Christianity, Israel, and the USA as a nation; he sees Dugin as calling for their extinction.


\hfill

\texttt{Cologero on 2012-09-06 at 00:02 said: }

No, Jason-Adam, I hadn't read it, despite its interesting premise. At the time of its publication, I would have been supporting the Left and would have seen no purpose in such an alliance. As far as its being ``first necessary to destroy the West", every death is a birth; what interregnum of the left born from that destruction would then cede power voluntarily to Tradition?

As Evola makes clear in Point 10, the bourgeoisie can be opposed in two directions: from below and from above. What purpose would be served by adopting the goals of the proletariat? Would ``aristocrats of the soul" take to the streets? We agree with Joseph de Maistre that what is necessary is not a revolution in the opposite direction, but rather the opposite of a revolution.

I'll reconsider it in the forthcoming review of Dugin's Fourth Political Theory.


\hfill

\texttt{Cologero on 2012-09-06 at 00:10 said: }

Michael, Carvalho himself anticipates a ``worldwide united front of Christians [Carvalho does not include Orthodox], Jews, and American nationalists", whose existence, he believes, is threatened internally by the globalist elite (Syndicate) and externally by the powers Dugin supports: BRIC and international Islam. It is certainly not obvious what would unite those three groups, other than the desire for mere survival. The Catholic Church, such as it is, seems unable, and even unwilling, to do anything to oppose the modern project. However, this is not to say that the Traditional spiritual forces of the West cannot arise again in some other way.


\hfill

\texttt{Jason-Adam on 2012-09-08 at 22:33 said: }

The end goal is to restore Tradition. The perspective Evola took later on in his life was that it was no longer possible to go BACK to the Tradition of our ancestors but that we need to move forward to the start of the new cycle, hence we support the proletariat to speed up the process of destruction.

Over a period of several generations, it is possible that a new aristocracy would emerge from communism, and the system would transform itself from a material system to a metaphysical based one…..the example of north Korea must be mentioned. north Korea has become a quite traditional monarchy based on Confucianism and (according to even its own officials) metaphysics.


\hfill

\texttt{Jason-Adam on 2012-09-08 at 22:39 said: }

The other point to mention is that, in his recent essay THE MISSION of JULIUS EVOLA, Dugin said that Guenon and Evola were wrong to declare capitalism the thought of the third estate, and communism that of the fourth. He makes an intruging thought that since the Kali Yuga concludes with the rise of the Counter Tradition, atheistic communism cannot be the Antichrist. The fact the USSR fell while the USA remains shows that in the world of decline the Soviet Union was not as fallen as America is.

I forget to add that these are not my views but I am arguing Dugin's case, as part of a Socratic exercise to stimulate Cologero and Logres.


\hfill

\texttt{Logres on 2012-09-08 at 22:56 said: }

Dugin wants to argue that the West is ``water" and hence Chaos (the Thalassocracy of the Sea). But isn't Water just one of the Four Elements? Why would (inherently) Water be worse than Earth (Asia)? And does Water have something to do with the spiritual flood God is pouring out over the Earth? And does this make the North ``Fire" and the South ``Air"? Dugin has some fascinating insights and criticisms of the Anglo imperium, most of which I agree with in point, but I'm waiting for C to review the book before saying more. The destructive Gnostic side of Dugin is evident here:

\url{http://arctogaia.com/public/eng/gnostic.htm}

Russia will always be ``Red" (Mouravieff emphasizes this) but I see no reason why Dugin would want to characterize the spiritual heritage of Orthodoxy as inherently ``Gnostic" in the sense he describes. Part of the problem is with terms, but another part of the problem is that different sections of the globe are going to be ``different" while remaining Traditional. Dugin may be deceived by some of his own rhetoric, it's possible – it wouldn't be the first time someone extremely bright (eg., Zizek) has gotten carried away by the sound of his own voice and ended up spinning webs of illusion for himself.


\hfill

\texttt{Cologero on 2012-09-09 at 18:45 said: }

I am going to ask you, Jason-Adam, to hold that thought until we begin a review of Dugin's book The Fourth Political Theory. Despite his vast and deep knowledge of Guenon and Evola, he makes some serious errors. The first and most obvious is confusing the empirical and contingent with principles. Guenon and Evola point to the degeneration of castes, not necessarily to specific ideologies. For example, the problem is not capitalism in itself; rather it is making the values of the bourgeois the dominant values, and so on.

Keep in mind, too, that the USSR was leninist; according to Marx, communism would arise only after capitalism, so the USSR was premature and does not reflect perfectly what Dugin tries to pull from it. Also, there is little doubt that proletarian values are predominating in the USA; this is what we need to ponder, not how we label its economic system. Beyond that, there is the ``fifth caste", the former outcastes, whose values (or lack thereof) predominate. How this manifests is up for discussion.

The worse fault is the metaphysics of Chaos over Logos, a very unfortunate mental error. The reasons should be obvious, but this will be a topic of discussion perhaps later this week.


\hfill

\texttt{Mihai on 2012-09-10 at 09:59 said: }

``it is an integral part of the Christian conception"

The problem with Evola, in respect to his view of Christianity, is that when he speaks unconditionally about ``Christianity" he actually has in mind Catholicism, which is certainly not the whole of Christianity. He is usually right when he criticizes certain failings of Catholicism, but- it is clear from his writings- that he completely ignores and has no knowledge of Orthodox (Eastern) Christianity.

It is true that here he only speaks for western Europe, but in general, he writes about Christianity as if there is nothing there except Catholicism. 

Also, I agree that when talking about a positive solution to bring Europe back to Tradition, he only gives equivocal solutions. 

``[….]for our men, pure reference to the spirit will suffice, precisely as the evidence of a transcendent reality, to appeal to another force in order to graft onto our force, to have a foreboding that our battle is not only a political battle, to attract an invisible consecration on a new world of men and leaders of men."

At best, this may only work for a few, very gifted, persons; but even they need a certain form in which to practice their spiritual aspirations. All in all, this really doesn't answer to the question of a ``practical" spiritual life.


\hfill

\texttt{Michael on 2012-09-10 at 10:49 said: }

Evola indicts Christianity as a whole: ``These forces worked spiritually, through primitive Christianity, to destroy the European spirit. At first, they concealed themselves within the lunar spirituality which took shape in the Catholic church, that is to say, a spirituality whose type is no longer the sacred king, the solar initiate, or the ``hero", but the saint or the priest who bows before God, whose ideal is no longer the warlike, sacral hierarchy and ``glory" but fraternal community and caritas."

I doubt that Orthodoxy is that different on these points given that the teachings of Catholicism and Orthodoxy both spring from the same Founder, the same scriptures. That is not to say that there are not differences, only that these differences are insubstantial in comparison to what is held in common.

I disagree with Evola on his rejection of Catholicism as the religious form for the West. The old imperial paganism and the cult of Mithra may be more ideal (from Evola's perspective), but they are both long dead because there is no succession in these religions.

I also don't follow on what has changed in the Catholic Church that has rendered it completely illegitimate. It was arguably valid in the Middle Ages. True, Western Civilization has changed, but the change has been in spite of the Church, not because of it.


\hfill

\texttt{Cologero on 2012-09-10 at 12:44 said: }

Evola has been remarkably inconsistent over the years on this question. This is probably one of the reasons for Guenon's accusation that Evola was full of prejudices. Orientations was updated toward the end of his life, as the references to Vatican II make clear. Nevertheless, if the church still supported the Syllabus (which I don't see how Evola could agree with in toto) and there were warrior orders, he would support Catholicism. Keep in mind that the purpose of those orders was to (1) destroy paganism and (2) recapture Jerusalem. Which of those programs does Evola support? Or who else should they fight? Perhaps liberals? Is open warfare really the best option today? These and other obvious questions are never addressed.

Guenon and Evola disagree fundamentally on this issue. Guenon asserts that Christianity was esoteric at its source, hence Western Medievalism was also. For him, there would be no problem with Orthodoxy, particularly since it preserved much of the neo-platonic tradition. For Evola, on the other hand, early christianity was an aberration that was corrected in the West only by its fortuitous encounter with Nordic influence. The East did not have that encounter, but I doubt Evola even gave it a moment's thought.

As for the church today, which is embarrassed by the Middle Ages and treats the Syllabus with neglect, how can it be the same? Vatican II opened the door to the modern world, ignoring completely both the crisis and the revolt. All that is necessary is to read, for example, De Monarchia or Consolations of Philosophy to realize there was quite a different thought process. Since the Middle Ages was the last Tradition in the West, it is the place to start. To mindlessly repeat Evola's prejudices, which we run into quite often, is to really miss the point.


\hfill

\texttt{Michael on 2012-09-10 at 14:34 said: }

Thanks for the clarification Cologero, but now my reading list has expanded yet again to include Boethius and the documents of the Council (read in light of the Syllabus).

What are your thoughts on the chivalric orders that have survived into the present? As you point out, they are no longer warrior organizations in our current age.


\hfill

\texttt{Mihai on 2012-09-10 at 15:48 said: }

The Eastern Roman Empire (known today as the Byzantine Empire), preserved the old Greco-Roman traditions and laws and incorporated them in a Christian synthesis. 

Not to mention the fact that Orthodoxy preserved throughout the centuries an initiatic path- Hesychia. Of course, this is a jnanic, contemplative path, which I doubt would attract much of Evola's interest. Still, I am pretty sure that he had no idea of it, because in one of his articles he said something like ``The notion of a Christian initiation is, today, a pure fantasy", while it is clear that Hesychia is not a fantasy, even though it is quite difficult to obtain such an initiation. 

I think Evola's view of Christianity was much too influenced by his early Nietzsche period, and this carried on, even if only residually, throughout the rest of his life.


\hfill

\texttt{Jason-Adam on 2012-09-11 at 01:47 said: }

Evola is very contradictory. Here it seems that if the beliefs of the Vatican hierarchy were the same as those of the SSPX, Evola would be pro-Catholic. Yet when the Church was under B. Pius XII (who is denigrated by the enemy az a nazi) and fighting communism with success in Spain and Portugal, not to forget Croatia, and working with fascism from above through the Society of Jesus to make it into a positive state-building force while eliminating the negative destructive elements much as Evola was trying to do himself, Evola was totally anti-Catholic that whole time !

For myself, I am trying to decide on whether to adopt a western Orthodox or traditional Catholic practice. I lean towards orthodoxy because it preserves the doctrine of imperium as well as initiation through hesychasm and deification/theosis; but the one flaw I find in orthodoxy is that it does not have a theology of holy war the way catholicism does. I actually think traditional Catholicism as being somewhat more masculine and warlike than orthodoxy 

Now, I think Evola's dismissal of orthodoxy comes from his reading DeMaistre's book on the Pope where he considers orthodoxy to be another form of protestantism.


\hfill

\texttt{Dominion on 2012-09-11 at 02:04 said: }

In place of the western Crusade doctrine, Orthodoxy appears to have its battle theology encapsulated in the histories of the Orthodox nations. Tsar Lazar of Serbia, who chose the Kingdom of Heaven over that of Earth and died in a holy battle, is an example.


\hfill

\texttt{Mihai on 2012-09-11 at 02:15 said: }

Catholicism seems to have created a confusion between temporal power and spiritual authority. The Pope's pretence to temporal power and the attempt at usurpation and annexation of the duties that rightfully belong to the Emperor have created many problems in the West.

In the East, there was no such confusion. Orthodoxy certainly is a path that is dedicated towards contemplation and purely metaphysical knowledge. However, this does not exclude the participation of more active types in the tradition.

The advantage that Orthodoxy has over Catholicism is that the latter has as its supreme point of reference the Supreme Being. The former goes beyond this point and considers the Supreme Identity and synthesis between Being and Non-Being. 

DeMaistre, while I greatly respect him, does have some limitations because of his exclusivist Catholic upbringing. Even his easy dismissal of Luther and early Protestantism doesn't stand, as the problem is more complex than he seemed to consider. I have seen this sectarian type of exclusivism in many Catholics, which is a really assumptive attitude on their part. After all, it was in Catholic Europe that the modern deviation began, and it was from there that it spread all over the world. Catholicism has not only been unable to do anything to stop it, but, with Vatican II, it even gave in to it and began accepting whatever terms modernism imposed on it.

As always, realization is indirectly proportional to the level of boasting.


\hfill

\texttt{Graham on 2012-09-11 at 11:58 said: }

Regarding Hesychastic initiation:

As we indicated above, genuine Hesychasm is so exclusive that any attempt at popularization and vulgarization is impossible.What is scattered in the modern world is a pseudo-Hesychasm which gives many the illusion that they obtain an initiation by repeating the Prayer of the Heart or practicing other techniques made public.

\url{http://www.scribd.com/doc/81326126/Hesychasm-II}


\hfill

\texttt{Graham on 2012-09-11 at 12:30 said: }

``Now, I think Evola's dismissal of orthodoxy comes from his reading DeMaistre's book on the Pope where he considers orthodoxy to be another form of protestantism."

Christ himself established the Petrine Primacy: `Thou art Peter and upon this rock I build my church…’ . This was accepted with docility by the Eastern Fathers for hundreds of years, until the wanton rebellion that led to the Great Schism. Under such conditions I'm not sure how one can believe that the East ``preserves the doctrine of Imperium."


\hfill

\texttt{Jason-Adam on 2012-09-11 at 19:30 said: }

I am willing to accept that, Graham, but do Christ's words extend to the temporal sphere of things ? Is Peter Caesar ?


\hfill

\texttt{Audrey on 2012-09-16 at 21:29 said: }

You know, I would agree 99\%. The missing percent being a caveat on the female gender.. Every female in the beginning of her sexual life is at the mercy of a good partner who will assist her to grow up. It rarely happens and predatorial female cougars abound.. not to mention gurus. Society would be less to the left if sexual perversity did not abound.. Yet males need to be held accountable and they never are. More often than not they like to play Lambs of God as if it is all the women's fault. I believe that Evola addressed that as well.


\hfill

\texttt{Andrew on 2012-09-17 at 23:01 said: }

Audrey is correct and that point deserves more discussion.

I have always thought a return to Sol Invictus is the only real option in a world that depends so much on the empirical. Manifestation is the living symbol of the principles from which phenomena precipitate. Emulation of the natural world's relation to the sun would produce a religion and society perfect on every level of realization.


\hfill

\texttt{Andrew on 2012-09-17 at 23:03 said: }

Yes, Evola did address that point to the tune of, ``the problem with women is that men have let them deviate in line with their natural tendencies".


\hfill

\texttt{Andrew on 2012-09-17 at 23:07 said: }

There is an entire class of mythology, of which the trials of Heracles are a part, that tell the story of the ``herculean" efforts it took to subjugate feminine nature to the order of Heaven. It was a sustained effort. What we have today is several generations that have given up on that civilizing project.


\end{sffamily}\end{footnotesize}
